\documentclass[twocolumn]{aastex631}
\begin{document}
\title{A residual reionization ionizing-photon-budget shortfall for star-forming galaxies at z\~{}9 (o32 systematic corner)}
\author{NebulaMind Lab (autonomous pipeline)}
\affiliation{NebulaMind Lab, \url{https://lab.nebulamind.net}}
\begin{abstract}
Reionization ionizing-photon-budget reconciliation at z\~{}9: star-forming galaxies require f\_esc=0.336 (+0.339/-0.172) to close the budget (Madau-Dickinson SFRD, log xi\_ion=25.5+/-0.15, clumping C=2-5, JWST-SFRD tail), versus indirect-proxy-inferred f\_esc=0.080 (+0.146/-0.051) from LzLCS O32/beta calibrations. Median delta(required-inferred)=+0.229 dex-frac (16-84\%: +0.029 to +0.568); 87\% of the systematic MC shows a shortfall. Conclusion: a genuine SHORTFALL remains, and the sign holds under both O32 and beta calibrations. The result is bounded by the xi\_ion x clumping x proxy-calibration systematic, not by statistics. Generated autonomously from public data (jwst) via the NebulaMind Lab runner: a bounded, reproducible, descriptive study.
\end{abstract}
\section{Introduction}
The reionization-photon-budget crisis has been a topic of interest in recent years, with studies suggesting that there may be a shortfall in the number of ionizing photons required to drive cosmic reionization [Muñoz2024]. This issue is further complicated by uncertainties in key parameters such as the escape fraction (f\_esc) and the ionizing efficiency (xi\_ion). Previous works have attempted to address this problem using various methods, including excursion set models [Park2022] and analytic approaches [Madau2017], but a comprehensive understanding remains elusive. The need for a more accurate assessment of the reionization-photon-budget is highlighted by the potential implications for our understanding of galaxy evolution and the early universe.
\section{Data and method}
To address this issue, we employ a literature-anchored budget calculation that does not rely on survey catalog data. Instead, we use the cosmic star formation rate density (SFRD) from Madau \& Dickinson's (2014) analytic fitting function as the foundation for our analysis. We adopt published values for xi\_ion and the O32/beta f\_esc proxy calibrations from LzLCS [Chisholm+22, Flury+22; Simmonds+24]. By combining these elements, we aim to reconcile the reionization ionizing-photon-budget at z\~{}9.
\section{Result}
Our analysis reveals that star-forming galaxies require an escape fraction of f\_esc=0.336 (+0.339/-0.172) to close the budget, based on the Madau-Dickinson SFRD and assuming log xi\_ion=25.5±0.15 with a clumping factor C between 2 and 5. However, indirect-proxy-inferred f\_esc values from LzLCS O32/beta calibrations suggest a significantly lower value of 0.080 (+0.146/-0.051). This discrepancy results in a median delta(required-inferred) of +0.229 dex-frac (16-84\%: +0.029 to +0.568), with 87\% of the systematic Monte Carlo simulations showing a shortfall.
\begin{figure}\centering\includegraphics[width=\columnwidth]{result.png}\caption{A residual reionization ionizing-photon-budget shortfall for star-forming galaxies at z\~{}9 (o32 systematic corner)}\end{figure}
\section{Caveats}
It is essential to acknowledge the limitations of our approach, which relies on automated, single-selection, and uncalibrated measurements. The accuracy of our results depends heavily on the assumptions made regarding xi\_ion, clumping factor C, and proxy calibrations. Additionally, our analysis does not account for potential systematic errors in the SFRD fitting function or uncertainties in the LzLCS data. These factors may introduce biases that affect the validity of our findings, highlighting the need for further research to refine our understanding of the reionization-photon-budget crisis.
\section*{References}
\footnotesize
[Muoz2024] Muñoz et al. (2024). Reionization after JWST: a photon budget crisis?. bibcode 2024MNRAS.535L..37M \par [Davies2021] Davies et al. (2021). The Predicament of Absorption-dominated Reionization: Increased Demands on Ionizing Sources. bibcode 2021ApJ...918L..35D \par [Park2022] Park et al. (2022). Calibrating excursion set reionization models to approximately conserve ionizing photons. bibcode 2022MNRAS.517..192P \par [Duncan2015] Duncan et al. (2015). Powering reionization: assessing the galaxy ionizing photon budget at z < 10. bibcode 2015MNRAS.451.2030D \par [Madau2017] Madau (2017). Cosmic Reionization after Planck and before JWST: An Analytic Approach. bibcode 2017ApJ...851...50M

\end{document}
